\documentclass[10pt,letterpaper]{article}
\usepackage[utf8]{inputenc}
\usepackage[intlimits]{amsmath}
\usepackage{amsfonts}
\usepackage{amssymb}
\usepackage{ragged2e}
\usepackage[letterpaper, margin=1in]{geometry}
\usepackage{graphicx}
\usepackage{cancel}
\usepackage{mathtools}
\usepackage{tabularx}
\usepackage{arydshln}
\usepackage{tensor}
\usepackage{array}
\usepackage{xcolor}
\usepackage[boxed]{algorithm}
\usepackage[noend]{algpseudocode}
\usepackage{listings}
\usepackage{textcomp}
\usepackage{mathrsfs}
\usepackage{bbm}
\usepackage{tikz}
\usepackage{enumitem}
\usepackage{arydshln}
\usepackage{relsize}
\usepackage{multicol}

\usetikzlibrary{bayesnet}

%%%%%%%%%%%%%%%%%%%%%%%%%%%%%
% Formatting commands
%%%%%%%%%%%%%%%%%%%%%%%%%%%%%
\newcommand{\n}{\hfill\break}
\newcommand{\up}{\vspace{-\baselineskip}}
\newcommand{\lemma}[1]{\par\noindent\settowidth{\hangindent}{\textbf{Lemma: }}\textbf{Lemma: }#1}
\newcommand{\defn}[1]{\par\noindent\settowidth{\hangindent}{\textbf{Defn: }}\textbf{Defn: }#1\n}
\newcommand{\thm}[1]{\par\noindent\settowidth{\hangindent}{\textbf{Thm: }}\textbf{Thm: }#1\n}
\newcommand{\prop}[1]{\par\noindent\settowidth{\hangindent}{\textbf{Prop: }}\textbf{Prop: }#1\n}
\newcommand{\cor}[1]{\par\noindent\settowidth{\hangindent}{\textbf{Cor: }}\textbf{Cor: }#1\n}
\newcommand{\ex}[1]{\par\noindent\settowidth{\hangindent}{\textbf{Ex: }}\textbf{Ex: }#1\n}
\newcommand{\exer}[1]{\par\noindent\settowidth{\hangindent}{\textbf{Exer: }}\textbf{Exer: }#1\n}
\newcommand{\proven}{\;$\square$\n}
\newcommand{\problem}[1]{\par\noindent{#1}\n}
\newcommand{\problempart}[2]{\par\noindent\indent{}\settowidth{\hangindent}{\textbf{(#1)} \indent{}}\textbf{(#1)} #2\n}
\newcommand{\ptxt}[1]{\textrm{\textnormal{#1}}}
\newcommand{\inlineeq}[1]{\centerline{$\displaystyle #1$}}
\newcommand{\pageline}{\noindent\rule{\textwidth}{0.1pt}}

%%%%%%%%%%%%%%%%%%%%%%%%%%%%%
% Math commands
%%%%%%%%%%%%%%%%%%%%%%%%%%%%%
% Set Theory
\newcommand{\card}[1]{\left|#1\right|}
\newcommand{\set}[1]{\left\{#1\right\}}
\newcommand{\setmid}{\;\middle|\;}
\newcommand{\ps}[1]{\mathcal{P}\left(#1\right)}
\newcommand{\pfinite}[1]{\mathcal{P}^{\ptxt{finite}}\left(#1\right)}
\newcommand{\naturals}{\mathbb{N}}
\newcommand{\N}{\naturals}
\newcommand{\integers}{\mathbb{Z}}
\newcommand{\Z}{\integers}
\newcommand{\rationals}{\mathbb{Q}}
\newcommand{\Q}{\rationals}
\newcommand{\reals}{\mathbb{R}}
\newcommand{\R}{\reals}
\newcommand{\complex}{\mathbb{C}}
\newcommand{\C}{\complex}
\newcommand{\halfPlane}{\mathbb{H}}
\let\H\relax
\newcommand{\H}{\halfPlane}
\newcommand{\comp}{^{\complement}}
\DeclareMathOperator{\Hom}{Hom}
\newcommand{\Ind}{\mathbbm{1}}
\newcommand{\cut}{\setminus}
\DeclareMathOperator{\elem}{elem}

% Graph Theory
\let\deg\relax
\DeclareMathOperator{\deg}{deg}
\newcommand{\degp}{\ptxt{deg}^{+}}
\newcommand{\degn}{\ptxt{deg}^{-}}
\newcommand{\precdot}{\mathrel{\ooalign{$\prec$\cr\hidewidth\hbox{$\cdot\mkern0.5mu$}\cr}}}
\newcommand{\succdot}{\mathrel{\ooalign{$\cdot\mkern0.5mu$\cr\hidewidth\hbox{$\succ$}\cr\phantom{$\succ$}}}}
\DeclareMathOperator{\cl}{cl}
\DeclareMathOperator{\affdim}{affdim}

% Probability
\newcommand{\parSymbol}{\P}
\newcommand{\Prob}{\mathbb{P}}
\renewcommand{\P}{\Prob}
\newcommand{\Avg}{\mathbb{E}}
\newcommand{\E}{\Avg}
\DeclareMathOperator{\Var}{Var}
\DeclareMathOperator{\cov}{cov}
\DeclareMathOperator{\Unif}{Unif}
\DeclareMathOperator{\Binom}{Binom}
\newcommand{\CI}{\mathrel{\text{\scalebox{1.07}{$\perp\mkern-10mu\perp$}}}}

% Standard Math
\newcommand{\inv}{^{-1}}
\newcommand{\abs}[1]{\left|#1\right|}
\newcommand{\ceil}[1]{\left\lceil{}#1\right\rceil{}}
\newcommand{\floor}[1]{\left\lfloor{}#1\right\rfloor{}}
\newcommand{\conj}[1]{\overline{#1}}
\newcommand{\of}{\circ}
\newcommand{\tri}{\triangle}
\newcommand{\inj}{\hookrightarrow}
\newcommand{\surj}{\twoheadrightarrow}
\newcommand{\ndiv}{\nmid}
\newcommand{\divides}{\mid}
\newcommand{\ndivides}{\nmid}
\DeclareMathOperator{\lcm}{lcm}
\DeclareMathOperator{\sgn}{sgn}
\newcommand{\map}[4]{\!\!\!\begin{array}[t]{rcl}#1 & \!\!\!\!\to & \!\!\!\!#2\\ #3 & \!\!\!\!\mapsto & \!\!\!\!#4\end{array}}
\newcommand{\bigsum}[2]{\smashoperator[lr]{\sum_{\scalebox{#1}{$#2$}}}}

% Linear Algebra
\newcommand{\Id}{\textrm{\textnormal{Id}}}
\newcommand{\im}{\textrm{\textnormal{im}}}
\newcommand{\norm}[1]{\abs{\abs{#1}}}
\newcommand{\tpose}{^{T}}
\newcommand{\iprod}[1]{\left<#1\right>}
\DeclareMathOperator{\trace}{tr}
\newcommand{\chgBasMat}[3]{\!\!\tensor*[_{#1}]{\left[#2\right]}{_{#3}}}
\newcommand{\vecBas}[2]{\tensor*[]{\left[#1\right]}{_{#2}}}
\DeclareMathOperator{\GL}{GL}
\DeclareMathOperator{\Mat}{Mat}
\DeclareMathOperator{\vspan}{span}
\DeclareMathOperator{\rank}{rank}
\newcommand{\V}[1]{\vec{#1}}
\DeclareMathOperator{\proj}{proj}
\DeclareMathOperator{\compProj}{comp}
\DeclareMathOperator{\row}{row}
\newcommand{\smallPMatrix}[1]{\paren{\begin{smallmatrix}#1\end{smallmatrix}}}
\newcommand{\smallBMatrix}[1]{\brack{\begin{smallmatrix}#1\end{smallmatrix}}}

% Multilinear Algebra
\newcommand{\Lsym}{\L}
\let\L\relax
\DeclareMathOperator{\L}{\mathscr{L}}
\DeclareMathOperator{\A}{\mathcal{A}}
\DeclareMathOperator{\Alt}{Alt}
\DeclareMathOperator{\Sym}{Sym}
\newcommand{\ot}{\otimes}
\newcommand{\ox}{\otimes}
\DeclareMathOperator{\asc}{asc}
\DeclareMathOperator{\asSet}{set}
\DeclareMathOperator{\sort}{sort}
\DeclareMathOperator{\ringA}{\mathring{A}}

% Topology
\newcommand{\closure}[1]{\overline{#1}}
\newcommand{\uball}{\mathcal{U}}
\DeclareMathOperator{\Int}{Int}
\DeclareMathOperator{\Ext}{Ext}
\DeclareMathOperator{\Bd}{Bd}
\DeclareMathOperator{\rInt}{rInt}
\DeclareMathOperator{\ch}{ch}
\DeclareMathOperator{\ah}{ah}
\newcommand{\LargerTau}{\mathlarger{\mathlarger{\mathlarger{\mathlarger{\tau}}}}}
\newcommand{\Tau}{\mathcal{T}}

% Analysis
\DeclareMathOperator{\Graph}{Graph}
\DeclareMathOperator{\epi}{epi}
\DeclareMathOperator{\hypo}{hypo}
\DeclareMathOperator{\supp}{supp}
\newcommand{\lint}[2]{\underset{#1}{\overset{#2}{{\color{black}\underline{{\color{white}\overline{{\color{black}\int}}\color{black}}}}}}}
\newcommand{\uint}[2]{\underset{#1}{\overset{#2}{{\color{white}\underline{{\color{black}\overline{{\color{black}\int}}\color{black}}}}}}}
\newcommand{\alignint}[2]{\underset{#1}{\overset{#2}{{\color{white}\underline{{\color{white}\overline{{\color{black}\int}}\color{black}}}}}}}
\newcommand{\extint}{\ptxt{ext}\int}
\newcommand{\extalignint}[2]{\ptxt{ext}\alignint{#1}{#2}}
\newcommand{\conv}{\ast}
\newcommand{\pd}[2]{\frac{\partial{}#1}{\partial{}#2}}
\newcommand{\del}{\nabla}
\DeclareMathOperator{\grad}{grad}
\DeclareMathOperator{\curl}{curl}
\let\div\relax
\DeclareMathOperator{\div}{div}
\DeclareMathOperator{\vol}{vol}

% Complex Analysis
\let\Re\relax
\DeclareMathOperator{\Re}{Re}
\let\Im\relax
\DeclareMathOperator{\Im}{Im}
\DeclareMathOperator{\Res}{Res}

% Proofs
\newcommand{\st}{s.t.}
\newcommand{\unique}{!}
\newcommand{\iffdef}{\overset{\ptxt{def}}{\Leftrightarrow}}
\newcommand{\eqdef}{\overset{\ptxt{def}}{=}}
\newcommand{\eqVertical}{\rotatebox[origin=c]{90}{=}}
\newcommand{\mapsfrom}{\mathrel{\reflectbox{\ensuremath{\mapsto}}}}
\newcommand{\mapsdown}{\rotatebox[origin=c]{-90}{$\mapsto$}\mkern2mu}
\newcommand{\mapsup}{\rotatebox[origin=c]{90}{$\mapsto$}\mkern2mu}
\newcommand{\from}{\!\mathrel{\reflectbox{\ensuremath{\to}}}}

% Brackets
\newcommand{\paren}[1]{\left(#1\right)}
\renewcommand{\brack}[1]{\left[#1\right]}
\renewcommand{\brace}[1]{\left\{#1\right\}}
\newcommand{\ang}[1]{\left<#1\right>}

% Algorithms
\algrenewcommand{\algorithmiccomment}[1]{\hskip 1em \texttt{// #1}}
\algrenewcommand\algorithmicrequire{\textbf{Input:}}
\algrenewcommand\algorithmicensure{\textbf{Output:}}
\newcommand{\algP}{\ptxt{\textbf{P}}}
\newcommand{\algNP}{\ptxt{\textbf{NP}}}
\newcommand{\algNPC}{\ptxt{\textbf{NP-Complete}}}
\newcommand{\algNPH}{\ptxt{\textbf{NP-Hard}}}
\newcommand{\algEXP}{\ptxt{\textbf{EXP}}}

%%%%%%%%%%%%%%%%%%%%%%%%%%%%%
% Other commands
%%%%%%%%%%%%%%%%%%%%%%%%%%%%%
\newcommand{\flag}[1]{\textbf{\textcolor{red}{#1}}}
\newcommand{\uSym}{\u}
\let\u\relax
\newcommand{\u}[1]{\underline{#1}}
\newcommand{\bSym}{\b}
\let\b\relax
\newcommand{\b}[1]{\textbf{#1}}
\newcommand{\iSym}{\i}
\let\i\relax
\newcommand{\i}[1]{\textit{#1}}
\newcommand{\rvline}{\hspace*{-\arraycolsep}\vline\hspace*{-\arraycolsep}}


%%%%%%%%%%%%%%%%%%%%%%%%%%%%%
% Make l's curvy in math environments
%%%%%%%%%%%%%%%%%%%%%%%%%%%%%
\mathcode`l="8000
\begingroup
\makeatletter
\lccode`\~=`\l
\DeclareMathSymbol{\lsb@l}{\mathalpha}{letters}{`l}
\lowercase{\gdef~{\ifnum\the\mathgroup=\m@ne \ell \else \lsb@l \fi}}%
\endgroup

\newcommand{\B}{
    \begin{tikzpicture}
    \filldraw [fill=red, draw=black] (0, 0) rectangle (0.37, 0.45);
    \draw [line width=0.5mm, white ] (0.1,0.08) -- (0.1,0.38);
    \draw[line width=0.5mm, white ] (0.1, 0.35) .. controls (0.2, 0.35) and (0.4, 0.2625) .. (0.1, 0.225);
    \draw[line width=0.5mm, white ] (0.1, 0.225) .. controls (0.2, 0.225) and (0.4, 0.1625) .. (0.1, 0.1);
    \end{tikzpicture}
}

\author{Thomas Cohn}
\title{Local Tangent Space Alignment Symbol Glossary}
\date{\today} % Can also use \today

\begin{document}
\maketitle
\setlength\RaggedRightParindent{\parindent}
\RaggedRight

\begin{tabularx}{\linewidth}{|l|>{\raggedright\arraybackslash}X|} \hline
	$\mathcal{F}$ & The data manifold. $\mathcal{F}$ is a parameterizable $d$-manifold in $\R^{m}$, with $d<m$.\\ \hline
	$d$ & The underlying dimension of the data manifold $\mathcal{F}$, i.e., the ``feature'' space.\\ \hline
	$m$ & The dimension of the higher dimensional ``input'' space.\\ \hline
	$f$ & $f:C\subset\R^{d}\to\R^{m}$ is a smooth map which parameterizes $\mathcal{F}$.\\ \hline
	$C$ & $C\subset\R^{d}$ is a compact subset of $\R^{d}$.\\ \hline
	$x_{i}$ & $x_{1},\ldots,x_{N}\in\R^{m}$ are the data points in the higher dimensional space.\\ \hline
	$\tau_{i}$ & $\tau_{1},\ldots,\tau_{N}\in\R^{d}$ are the corresponding feature space coordinates of the $x_{i}$'s, where $x_{i}=f(\tau_{i})+\epsilon_{i}$.\\ \hline
	$\epsilon_{i}$ & $\epsilon_{1},\ldots,\epsilon_{N}\in\R^{d}$ is the added noise to $f(\tau_{i})$ to obtain $x_{i}$.\\ \hline
	$X$ & $X=\smallBMatrix{| & & |\\ x_{1} & \cdots & x_{N}\\ | & & |}$ is the matrix form of the data.\\ \hline
	$J_{f}$ & $J_{f}(\tau)=\smallBMatrix{\partial{}f_{1}/\partial\tau_{1} & \cdots & \partial{}f_{1}/\partial\tau_{d}\\ \vdots & \ddots & \vdots\\ \partial{}f_{m}/\partial\tau_{1} & \cdots & \partial{}f_{m}/\partial\tau_{d}}$ is the Jacobi matrix of $f$ at $\tau$.\\ \hline
	$\overline{\tau}$ & $\overline{\tau}$ is some point near $\tau$ used for the following approximation via Taylor expansion: $f(\overline{\tau})=f(\tau)+J_{f}(\tau)\cdot(\overline{\tau}-\tau)+O(\norm{\overline{\tau}-\tau}^{2})$.\\ \hline
	$\Tau_{\tau}$ & $\Tau_{\tau}=\vspan(J_{f}(\tau))$ is the tangent space of $f$ at $\tau$.\\ \hline
	$Q_{\tau}$ & $Q_{\tau}$ is a matrix forming some orthonormal basis of $\Tau_{\tau}$, allowing us to write $J_{f}(\tau)(\overline{\tau}-\tau)=Q_{\tau}\theta_{\tau}^{*}$.\\ \hline
	$\theta_{\tau}^{*}$ & ??? Possibly defined by $Q_{\tau}\theta_{\tau}^{*}=f(\overline{\tau})-f(\tau)$?\\ \hline
	$P_{\tau}$ & Defined by $\theta_{\tau}^{*}=Q_{\tau}\tpose{}J_{f}(\tau)(\tau-\overline{\tau})\equiv{}P_{\tau}(\overline{\tau}-\tau)$.\\ \hline
	$\theta_{\tau}$ & The mapping from $\tau$ to $\theta^{*}_{\tau}$ is a local affine transformation. $\theta_{\tau}$ is the orthogonal projection of $f(\overline{\tau})-f(\tau)$ onto $\Tau_{\tau}$. $\theta_{\tau}\equiv{}Q_{\tau}\tpose(f(\overline{\tau})-f(\tau))=\theta_{\tau}^{*}+O(\norm{\overline{\tau}-\tau}^{2})\approx\theta_{\tau}^{*}$.\\ \hline
	$\Omega(\tau)$ & $\Omega(\tau)$ defines the neighborhood of $\tau$.\\ \hline
	$X_{i}$ & $X_{i}=\smallBMatrix{| & & |\\ x_{i_{1}} & \cdots & x_{i_{k}}\\ | & & |}$ is a matrix consisting of the $k$-nearest neighbors of $x_{i}$.\\ \hline
	$x$ & $x$ is part of the minimization problem\n
	$\displaystyle\underset{x,\Theta,Q}{\min}\sum_{j=1}^{k}\norm{x_{i_{j}}-(x+Q\theta_{j})}_{2}^{2}=\underset{x,\Theta,Q}{\min}\norm{X_{i}-(xe\tpose+Q\Theta)}_{2}^{2}$.\n
	The optimal $x$ is $\overline{x}_{i}$ -- the mean of all the $x_{i_{j}}$'s.\\ \hline
	$Q$ & $Q$ is another term in the minimzation (but what is it???). It is of $d$ columns. Optimal $Q$ is given by the $d$ singular vectors of $X_{i}(I-ee\tpose/k)$ corresponding with the $d$ largest singular values.\\ \hline
	$\Theta$ & $\Theta=\smallBMatrix{| & & |\\ \theta_{1} & \cdots & \theta_{k}\\ | & & |}$ is a matrix consisting of the orthogonal projections of each $x_{i_{j}}-x$ onto $\Tau_{x_{i}}$. $\Theta$ is given by $\Theta_{i}=Q_{i}\tpose{}X_{i}(I-\frac{1}{k}ee\tpose)=\smallBMatrix{| & & |\\ \theta_{1}^{(i)} & \cdots & \theta_{k}^{(i)}\\ | & & |},\theta_{j}^{(i)}=Q_{i}\tpose(x_{i_{j}}-\overline{x}_{i})$.\\ \hline
	$\xi_{j}^{(i)}$ & $\xi_{j}^{(i)}=(I-Q_{i}Q_{i}\tpose)(x_{i_{j}}-\overline{x}_{i})$ is the global reconstruction error of $x_{i_{j}}$, with $x_{i_{j}}=\overline{x}_{i}+Q_{i}\theta_{j}^{(i)}+\xi_{j}^{(i)}$.\\ \hline
	$\tau_{i_{j}}$ & $\tau_{i_{j}}$ is the feature space coordinate of $x_{i_{j}}$. We want $\tau_{i_{j}}=\overline{\tau_{i}}+L_{i}\theta_{j}^{(i)}+\epsilon_{j}^{(i)}$, for $j\in\set{1,\ldots,k}$ and $i\in\set{1,\ldots,N}$.\\ \hline
	$\overline{\tau_{i}}$ & $\overline{\tau_{i}}$ is the mean of the $\tau_{i_{j}}$, for $j\in\set{1,\ldots,k}$.\\ \hline
	$L_{i}$ & $L_{i}$ is the local affine transformation for $\tau_{i}$.\\ \hline
	$E_{i}$ & $E_{i}=\smallBMatrix{| & & |\\ \epsilon_{1}^{(i)} & \cdots & \epsilon_{k}^{(i)}\\ | & & |}$ is the matrix form of the local reconstruction error.\\ \hline
	$e$ & $e=\smallBMatrix{1 & \cdots & 1}\tpose$ is a column vector of all ones of dimension $k$.\\ \hline
	$T_{i}$ & $T_{i}=\smallBMatrix{| & & |\\ \tau_{i} & \cdots & \tau_{k}\\ | & & |}$ is the matrix form of the feature space coordinates. It satisfies $T_{i}=\frac{1}{k}T_{i}ee\tpose+L_{i}\Theta_{i}+E_{i}$, so $E_{i}=T_{i}(I_{k}-\frac{1}{k}ee\tpose)-L_{i}\Theta_{i}$.\\ \hline
	$S_{i}$ & $S_{i}$ is the $0$-$1$ selection matrix such that $TS_{i}=T_{i}$.\\ \hline
	$S$ & $S=\left[\begin{array}{c;{2pt/2pt}c;{2pt/2pt}c}S_{1} & \cdots & S_{N}\end{array}\right]$.\\ \hline
	$W_{i}$ & $W_{i}=(I_{k}-\frac{1}{k}ee\tpose)(I-\Theta_{i}^{+}\Theta_{i})$.\\ \hline
\end{tabularx}
\begin{tabularx}{\linewidth}{|l|>{\raggedright\arraybackslash}X|} \hline
	$W$ & $W=\operatorname{diag}(W_{1},\ldots,W_{N})=\left[\begin{array}{c;{2pt/2pt}c;{2pt/2pt}c;{2pt/2pt}c}W_{1} & 0 & \cdots & 0\\ \hdashline[2pt/2pt] 0 & W_{2} & \cdots & 0\\ \hdashline[2pt/2pt] \vdots & \vdots & \ddots & \vdots\\ \hdashline[2pt/2pt] 0 & 0 & \cdots & W_{N}\end{array}\right]$.\\ \hline
	$T$ & We want to find $T$ to minimize the overall reconstruction error $\displaystyle\sum_{i}\norm{E_{i}}_{F}^{2}=\norm{TSW}_{F}^{2}$.\n
	To determine $T$, we require $TT\tpose=I_{d}$.\\ \hline
	$B$ & $B\equiv{}SWW\tpose{}S\tpose$ has eigenvector $e$ corresponding to eigenvalue $0$. Thus, optimal $T$ is given by the $d$ eigenvectors corresponding to the $2$nd to $d+1$st smallest eigenvalues.\\ \hline
\end{tabularx}

\end{document}